\chapter{Getting started}
\label{ch:getting-started}

Immersive Engineering does not begin with a machine. It begins with a hammer and a fire.

\section{The Engineer's Hammer}

Three iron ingots and two sticks. It is the mod's universal tool and you will never stop carrying
one:

\begin{itemize}
\item \textbf{Assembles multiblocks.} Every structure in \autoref{ch:multiblocks} is built by
      placing its blocks and then striking one of them.
\item \textbf{Rotates blocks.} Right-click most machines to turn them.
\item \textbf{Configures sides.} Sneak-right-click a machine face to cycle what that side does.
\item \textbf{Dismantles.} Takes a multiblock apart into its parts.
\end{itemize}

\section{The Engineer's Manual}

One book and one lever. It is the in-game version of this document, and it updates itself as you
obtain things. \autoref{ch:ingame} reproduces it in full.

\section{Treated wood}

Creosote oil in a barrel or a Bottling Machine turns planks into \textbf{treated wood}, which
resists weather and is the base of almost everything early. You get creosote from the coke oven, so
the coke oven comes first.

\section{The Coke Oven}

A $3\times3\times3$ of Coke Bricks. Turns coal into \textbf{coal coke} and \textbf{creosote oil},
both of which you need continuously:

\begin{itemize}
\item Coal coke burns longer than coal and is the fuel for the Blast Furnace.
\item Creosote oil treats wood and, later, is a Squeezer input.
\end{itemize}

\begin{ietip}
Build two. The oven is slow, both outputs are consumed constantly, and the second one costs almost
nothing next to the time it saves.
\end{ietip}

\section{The Blast Furnace}

A $3\times3\times3$ of Blast Bricks. Iron plus coal coke gives \textbf{steel}, which is the material
the entire middle of the mod is built from. It also produces \textbf{slag}, a Blast Brick
ingredient, so the furnace partly pays for its own expansion.

The \textbf{Advanced Blast Furnace} is the upgrade: faster, and it accepts external heating from a
Blast Furnace Preheater.

\section{Your first power}

Immersive Engineering measures power in \textbf{Immersive Flux (IF)}, which is Forge Energy under
another name --- it interoperates with RF and FE at one to one.

The cheapest start is a \textbf{Water Wheel} or a \textbf{Windmill}: both are built from treated
wood, need no fuel and produce a trickle. Put a \textbf{Dynamo} behind the wheel, an \textbf{LV
Connector} on the dynamo, and run copper wire to whatever needs it.

\begin{iewarning}
Wires are not decorative. Walking into an energised HV wire will hurt you, and connecting a wire to
a connector of the wrong tier will destroy it. See \autoref{ch:wiring}.
\end{iewarning}

\section{The order most people build in}

\begin{enumerate}
\item Hammer, manual, coke oven.
\item Treated wood, then a water wheel or windmill for a first trickle of power.
\item Blast furnace, then steel.
\item Crusher, for doubled ore.
\item Metal Press, for cheap plates and rods --- this is the point the mod starts paying you back.
\item Diesel generator or one of the fork's power plants (\autoref{ch:petroleum}) for real output.
\end{enumerate}
